% 图 72.2 —— 由 `checks/emit_hand.py` 自动拼出（probe 精测 + prims2 图元分解）
%%
% ⚠ 这是**稿子不是成品**：跑 `checks/checkhand.py 72.2` 看指标 + 看叠合图。
%%
%% ⚠ 待核：
%   · 本图 probe 报出阴影族 1 个 —— **本脚本不生成阴影**（要照原书手写）；prims2 的 strokes 里可能含阴影线，核对时留意别画重
%   · 可疑字形 7 项（空心圈 ⊙ 常被认成字母 o）—— 见文件末
%%
\documentclass[12pt]{standalone}
\usepackage{fontspec}
\setsansfont{Arial}
\newfontfamily\mfallback{DejaVu Sans}
\usepackage{tikz}
\begin{document}
\begin{tikzpicture}[x=1pt, y=-1pt, line width=4.0pt, line cap=round, line join=round]

  %% ════ 填充区（prims2 的 fills）════
  \fill (131.0,72.0) -- (145.0,88.0) -- (147.0,84.0) -- (135.0,72.0) -- cycle;

  %% ════ 直线段（prims2 的 strokes）════
  \draw[line width=5.7pt] (328.0,137.0) -- (326.0,142.0);
  \draw[line width=5.0pt] (361.0,212.0) -- (337.0,236.0);
  \draw[line width=5.0pt] (398.4,170.6) -- (403.0,167.0);
  \draw[line width=6.0pt] (403.0,167.0) -- (403.0,161.0);
  \draw[line width=3.0pt] (452.0,263.0) -- (452.0,260.0);
  \draw[line width=4.0pt] (452.0,263.0) -- (428.0,279.0);
  \draw[line width=4.0pt] (279.0,264.0) -- (307.0,254.0);
  \draw[line width=4.0pt] (340.0,359.0) -- (327.0,372.0);
  \draw[line width=3.0pt] (367.0,293.0) -- (370.0,293.0);
  \draw[line width=5.0pt] (367.0,293.0) -- (337.0,302.0);
  \draw[line width=5.0pt] (408.0,182.0) -- (415.0,192.0);
  \draw[line width=8.0pt] (459.0,253.0) -- (453.0,259.0);
  \draw[line width=5.0pt] (420.0,206.0) -- (429.0,220.0);
  \draw[line width=7.0pt] (371.0,292.0) -- (374.0,286.0);
  \draw[line width=6.5pt] (326.0,144.0) -- (334.0,155.0);
  \draw[line width=4.0pt] (505.0,281.0) -- (495.0,277.0);
  \draw[line width=4.0pt] (402.0,287.0) -- (426.0,279.0);
  \draw[line width=5.0pt] (100.0,157.0) -- (82.0,150.0);
  \draw[line width=5.0pt] (132.0,72.0) -- (145.0,85.0);
  \draw[line width=6.0pt] (291.0,347.0) -- (267.0,305.0);
  \draw[line width=3.5pt] (452.0,259.0) -- (449.0,246.0);
  \draw[line width=5.3pt] (320.0,145.0) -- (325.0,144.0);
  \draw[line width=4.0pt] (149.0,130.0) -- (137.0,147.0);
  \draw[line width=6.1pt] (267.0,304.0) -- (271.0,299.5);
  \draw[line width=4.0pt] (271.0,299.5) -- (279.0,265.0);
  \draw[line width=8.9pt] (266.0,305.0) -- (261.0,312.0);
  \draw[line width=4.0pt] (364.0,335.0) -- (341.0,358.0);
  \draw[line width=4.0pt] (482.0,270.0) -- (461.0,253.0);
  \draw[line width=7.6pt] (142.0,22.0) -- (148.0,18.0);
  \draw[line width=5.0pt] (305.0,327.0) -- (292.0,347.0);
  \draw[line width=7.3pt] (153.0,111.0) -- (152.0,118.0);
  \draw[line width=5.0pt] (426.0,310.0) -- (451.0,303.0);
  \draw[line width=6.3pt] (84.0,66.0) -- (85.0,72.0);
  \draw[line width=5.0pt] (452.0,302.0) -- (480.0,295.0);
  \draw[line width=4.0pt] (489.0,260.0) -- (506.0,280.0);
  \draw[line width=6.5pt] (460.0,252.0) -- (450.0,245.0);
  \draw[line width=4.0pt] (396.0,320.0) -- (365.0,334.0);
  \draw[line width=7.6pt] (429.0,221.0) -- (427.0,228.0);
  \draw[line width=4.0pt] (336.0,303.0) -- (307.0,321.0);
  \draw[line width=5.0pt] (258.0,269.0) -- (266.0,304.0);
  \draw[line width=5.0pt] (153.0,109.0) -- (152.0,98.0);
  \draw[line width=3.8pt] (520.0,119.0) -- (517.8,122.0);
  \draw[line width=4.5pt] (154.0,110.0) -- (168.0,111.0);
  \draw[line width=7.0pt] (72.0,81.0) -- (85.0,72.0);
  \draw[line width=5.0pt] (171.0,111.0) -- (200.0,111.0);
  \draw[line width=4.0pt] (62.0,110.0) -- (68.0,90.0);
  \draw[line width=5.0pt] (373.0,293.0) -- (398.0,287.0);
  \draw[line width=4.0pt] (306.0,326.0) -- (306.0,322.0);
  \draw[line width=4.0pt] (278.0,264.0) -- (259.0,268.0);
  \draw[line width=5.0pt] (257.0,268.0) -- (218.5,274.0);
  \draw[line width=5.0pt] (57.3,491.3) -- (77.0,509.0);
  \draw[line width=4.0pt] (77.0,509.0) -- (108.5,520.0);
  \draw[line width=4.0pt] (309.0,396.0) -- (326.0,373.0);
  \draw[line width=6.6pt] (92.0,68.0) -- (87.0,72.0);
  \draw[line width=4.0pt] (397.0,320.0) -- (425.0,310.0);
  \draw[line width=4.0pt] (308.0,254.0) -- (336.0,236.0);

  %% ════ 曲线（prims2 的 curves —— 三段式，坐标同为 (y,x)）════
  \draw[line width=4.0pt] (449.0,122.0) .. controls (433.3,142.7) and (422.1,167.0) .. (416.0,192.0);
  \draw[line width=5.0pt] (492.0,111.0) .. controls (502.1,110.7) and (512.1,111.9) .. (520.0,118.0);
  \draw[line width=4.0pt] (386.0,240.0) .. controls (382.8,215.8) and (389.3,191.9) .. (398.4,170.6);
  \draw[line width=4.0pt] (452.0,263.0) .. controls (458.7,275.6) and (468.3,285.6) .. (480.0,294.0);
  \draw[line width=4.0pt] (514.0,161.0) .. controls (510.6,155.2) and (516.6,149.1) .. (523.0,151.0);
  \draw[line width=4.0pt] (367.0,293.0) .. controls (351.3,269.6) and (356.8,239.1) .. (362.0,213.0);
  \draw[line width=5.0pt] (515.0,162.0) .. controls (521.6,163.6) and (525.8,158.2) .. (524.0,152.0);
  \draw[line width=4.0pt] (514.0,163.0) .. controls (497.3,197.0) and (502.9,246.4) .. (538.0,267.0);
  \draw[line width=4.0pt] (444.0,211.0) .. controls (447.1,220.9) and (441.4,234.8) .. (449.0,244.0);
  \draw[line width=4.0pt] (558.0,243.0) .. controls (516.3,230.6) and (520.2,165.6) .. (559.0,152.0);
  \draw[line width=4.0pt] (558.0,243.0) .. controls (573.4,217.4) and (577.4,178.5) .. (561.0,152.0);
  \draw[line width=4.0pt] (558.0,243.0) .. controls (555.0,252.5) and (547.3,260.9) .. (539.0,267.0);
  \draw[line width=5.0pt] (337.0,237.0) .. controls (331.0,256.6) and (326.8,281.4) .. (336.0,301.0);
  \draw[line width=4.0pt] (428.0,279.0) .. controls (432.6,289.3) and (442.3,295.2) .. (451.0,301.0);
  \draw[line width=5.0pt] (521.0,119.0) .. controls (529.7,120.3) and (540.0,125.0) .. (544.0,133.0);
  \draw[line width=5.0pt] (448.0,121.0) .. controls (430.0,127.9) and (414.1,142.5) .. (403.0,159.0);
  \draw[line width=4.0pt] (308.0,255.0) .. controls (303.5,276.0) and (294.1,300.6) .. (306.0,321.0);
  \draw[line width=4.0pt] (481.0,295.0) .. controls (490.2,293.2) and (503.5,293.1) .. (508.0,284.0);
  \draw[line width=4.0pt] (452.0,168.0) .. controls (459.2,146.3) and (472.5,126.4) .. (491.0,112.0);
  \draw[line width=5.0pt] (402.0,160.0) .. controls (383.7,170.8) and (372.3,193.0) .. (362.0,211.0);
  \draw[line width=4.0pt] (450.0,121.0) .. controls (460.9,111.4) and (476.8,108.1) .. (491.0,111.0);
  \draw[line width=5.0pt] (246.0,338.0) .. controls (233.1,408.4) and (149.6,460.4) .. (81.0,434.0);
  \draw[line width=3.0pt] (426.0,278.0) .. controls (413.0,251.7) and (412.8,222.3) .. (416.0,193.0);
  \draw[line width=4.0pt] (524.0,150.0) .. controls (527.2,142.2) and (536.1,137.7) .. (543.0,134.0);
  \draw[line width=4.0pt] (326.0,372.0) .. controls (313.5,367.1) and (300.3,359.5) .. (292.0,348.0);
  \draw[line width=3.0pt] (336.0,303.0) .. controls (340.6,316.7) and (351.0,327.5) .. (364.0,334.0);
  \draw[line width=4.0pt] (517.8,122.0) .. controls (483.2,139.2) and (476.1,180.9) .. (475.0,216.0);
  \draw[line width=3.0pt] (372.0,294.0) .. controls (371.1,308.2) and (386.5,312.7) .. (396.0,319.0);
  \draw[line width=5.0pt] (509.0,283.0) .. controls (519.6,284.6) and (531.3,276.5) .. (538.0,268.0);
  \draw[line width=3.0pt] (241.0,102.0) .. controls (272.7,104.2) and (302.6,117.9) .. (325.0,142.0);
  \draw[line width=4.0pt] (340.0,358.0) .. controls (325.8,352.1) and (313.7,340.5) .. (306.0,327.0);
  \draw[line width=5.0pt] (218.5,274.0) .. controls (149.4,275.0) and (80.7,306.0) .. (42.0,366.0);
  \draw[line width=5.0pt] (42.0,366.0) .. controls (24.7,405.7) and (30.8,457.1) .. (57.3,491.3);
  \draw[line width=5.0pt] (108.5,520.0) .. controls (193.5,528.5) and (269.5,467.0) .. (309.0,396.0);
  \draw[line width=4.0pt] (425.0,309.0) .. controls (416.5,304.0) and (405.6,298.3) .. (402.0,288.0);
  \draw[line width=5.0pt] (545.0,134.0) .. controls (552.8,135.3) and (557.4,144.5) .. (560.0,151.0);
  \draw[line width=5.0pt] (230.0,352.0) .. controls (178.4,354.1) and (127.0,382.1) .. (100.0,427.0);
  \draw[line width=4.0pt] (64.0,70.0) .. controls (61.2,66.3) and (68.1,58.2) .. (71.0,65.0);
  \draw[line width=4.0pt] (71.0,65.0) .. controls (73.5,71.5) and (66.1,79.0) .. (64.0,70.0);

  %% ════ 圆 / 椭圆（probe 精测）════
  \draw (108.1,111.5) circle (93.6);

  %% ════ 实心点 ════
  \fill (170.4,106.5) circle (4.6);
  \fill (171.5,116.0) circle (5.6);
  \fill (463.0,249.5) circle (6.1);
  \fill (274.0,307.5) circle (4.6);
  \fill (298.5,349.0) circle (4.0);
  \fill (291.5,352.5) circle (7.1);

  %% ════ 标签（probe 直接给的 \node 行）════
  %% ⚠ 全部补了 fill=white：文字在最上层，否则与曲线交叉干扰
     \node[anchor=center,inner sep=0pt,fill=white,font=\fontsize{29.1}{36.4}\selectfont] at (30.0,30.0) {\textsf{\textbf{\textit{f}}}};   % box(24,18,11,23)
     \node[anchor=center,inner sep=0pt,fill=white,font=\fontsize{27.1}{33.9}\selectfont] at (60.5,56.5) {\textsf{\textbf{\textit{f}}}};   % box(55,45,10,22)
     \node[anchor=center,inner sep=0pt,fill=white,font=\fontsize{44.6}{55.8}\selectfont] at (123.3,73.2) {{\mfallback\textbf{−}}};   % box(102,63,21,7)
     \node[anchor=center,inner sep=0pt,fill=white,font=\fontsize{29.6}{37.0}\selectfont] at (174.0,83.0) {\textsf{\textbf{\textit{Y}}}};   % box(165,70,17,24)
     \node[anchor=center,inner sep=0pt,fill=white,font=\fontsize{28.4}{35.5}\selectfont] at (302.2,91.7) {\textsf{\textbf{\textit{b}}}};   % box(293,80,16,21)
     \node[anchor=center,inner sep=0pt,fill=white,font=\fontsize{28.6}{35.8}\selectfont] at (109.2,113.4) {\textsf{\textbf{\textit{o}}}};   % box(101,104,15,16)
     \node[anchor=center,inner sep=0pt,fill=white,font=\fontsize{29.2}{36.4}\selectfont] at (220.9,129.0) {\textsf{\textbf{\textit{P}}}};   % box(210,117,18,22)
     \node[anchor=center,inner sep=0pt,fill=white,font=\fontsize{61.1}{76.4}\selectfont] at (68.0,147.3) {\textsf{\textbf{′}}};   % box(63,122,12,20)
     \node[anchor=center,inner sep=0pt,fill=white,font=\fontsize{30.5}{38.1}\selectfont] at (162.5,140.8) {\textsf{\textbf{\textit{q}}}};   % box(153,128,17,22)
     \node[anchor=center,inner sep=0pt,fill=white,font=\fontsize{47.7}{59.6}\selectfont] at (130.0,164.2) {{\mfallback\textbf{−}}};   % box(108,153,21,8)
     \node[anchor=center,inner sep=0pt,fill=white,font=\fontsize{34.0}{42.5}\selectfont] at (439.0,187.9) {\textsf{\textbf{9}}};   % box(428,176,18,22)
     \node[anchor=center,inner sep=0pt,fill=white,font=\fontsize{24.1}{30.2}\selectfont] at (453.0,196.2) {\textsf{\textbf{0}}};   % box(445,188,12,16)
     \node[anchor=center,inner sep=0pt,fill=white,font=\fontsize{30.7}{38.3}\selectfont] at (485.8,239.8) {\textsf{\textbf{\textit{b}}}};   % box(476,227,17,23)
     \node[anchor=center,inner sep=0pt,fill=white,font=\fontsize{33.8}{42.2}\selectfont] at (387.3,264.8) {\textsf{\textbf{6}}};   % box(379,251,15,26)
     \node[anchor=center,inner sep=0pt,fill=white,font=\fontsize{34.8}{43.5}\selectfont] at (236.0,312.4) {\textsf{\textbf{9}}};   % box(225,300,18,23)
     \node[anchor=center,inner sep=0pt,fill=white,font=\fontsize{33.0}{41.2}\selectfont] at (286.6,375.1) {\textsf{\textbf{\textit{a}}}};   % box(277,365,17,17)

  % 钉死画布：否则超出的图元/控制点会把页面撑大，1:1 回比就失准
  \pgfresetboundingbox
  \path[use as bounding box] (0,0) rectangle (584,533);
\end{tikzpicture}
\end{document}

% ⚠ 待核清单（probe 拿不准，人眼过一遍）：
%   · 可疑字形：\node[anchor=center,inner sep=0pt,font=\fontsize{28.6}{35.8}\selectfont] at (109.2,113.4) {\textsf{\textbf{\te
%   · 未识别/跳过：⑤ 跳过/未识别的字形
%   · 未识别/跳过：box(72,63,22,20)
%   · 未识别/跳过：box(131,72,17,17)
%   · 未识别/跳过：box(136,130,15,19)
%   · 未识别/跳过：box(81,148,21,11)
%   · 未识别/跳过：box(420,206,20,24)
